\begin{figure}[htbp]
\centering
\begin{tikzpicture}[x=.98cm,y=.98cm,>=Latex,font=\small,
  region/.style={draw=black!55,thick},
  walk/.style={draw=blue!65!black,very thick},
  cut/.style={draw=black!60,densely dashed}]
  \coordinate (a) at (0,0);
  \coordinate (r) at (2.30940,0);
  \coordinate (p1) at (1.27017,1.8);
  \coordinate (p2) at (3.69504,2.4);
  \coordinate (b) at (6.3,0);

  \path[region,fill=blue!7] (-2.30940,0)--(r)--(0,4)--cycle;
  \path[region,fill=orange!13] (r)--(.23094,3.6)--(4.38786,3.6)--cycle;
  \path[region,fill=green!7] (r)--(7.85196,0)--(5.08068,4.8)--cycle;
  \draw[thick] (-2.65,0)--(8.2,0);
  \node[anchor=north] at (4.5,-.60) {original base: $n_0\cdot(z-a)=0$};

  \draw[cut] (2.54034,-.4)--(-.80829,5.4)
    node[above] {$n_1\cdot(z-a)=t$};
  \draw[cut] (2.07846,-.4)--(5.42709,5.4)
    node[above] {$n_2\cdot(z-a)=t$};

  \draw[<->,black!55,thin] (1.27017,.03)--(1.27017,1.77)
    node[midway,right=3pt] {$q$};
  \draw[<->,black!55,thin] (3.69504,.03)--(3.69504,2.37)
    node[midway,left=3pt] {$X$};

  \draw[walk,->] (a) .. controls (0,.65) and (-.18,1.02) .. (.35,1.35)
    .. controls (.68,1.52) and (.94,1.61) .. (p1);
  \draw[walk,->] (p1) .. controls (1.83,2.12) and (2.01,2.61) .. (2.52,2.58)
    .. controls (3.03,2.55) and (3.26,2.65) .. (p2);
  \draw[walk,->] (p2) .. controls (4.18,2.12) and (4.37,2.55) .. (4.9,2.44)
    .. controls (5.8,2.15) and (6.3,.8) .. (b);

  \node[align=center] at (-.2,1.9) {$T(a,n_0,t)$};
  \node[align=center] at (2.16,3.22) {$T(p_1,n_1,q)$};
  \node[align=center] at (5.08,2.9) {$T(p_2,n_2,X)$};
  \fill (a) circle (1.8pt) node[below left=2pt] {$a$};
  \fill (p1) circle (1.8pt) node[above left=2pt] {$p_1$};
  \fill (p2) circle (1.8pt) node[above left=3pt] {$p_2$};
  \fill (b) circle (1.8pt) node[below right=2pt] {$b$};
  \draw[->] (-2.35,3.65)--(-2.35,4.4) node[above] {$n_0$};
  \draw[->] (-2.35,3.65)--(-1.7,4.025) node[right] {$n_1$};
  \draw[->] (-2.35,3.65)--(-1.7,3.275) node[right] {$n_2$};
\end{tikzpicture}
\caption{The three exits producing a half-plane return from $a$ to $b$.
The depths are $q=n_0\cdot(p_1-a)$ and $X=n_0\cdot(p_2-a)$.
The first triangle lies before the first dashed line, the second lies
after that line and before the second, and the third lies after the
second line and above the original base. Thus their interiors are
disjoint, and each dashed line is crossed once by the concatenation.
The triangle boundaries have the indicated directions; the blue paths
are schematic and the lattice is omitted.}
\label{fig:geometry-three-exits}
\end{figure}
